\begin{proposition}[The two complex beta integrals]\hypertarget{tag-001U}{}\label{ari:C.8.7}
\StatementRecord{ari:C.8.7}{proved}{001U}{\SeriesRef{ari:C.8.6}, \SeriesRef{ari:eq:ca-mahler}}

For the square and equianharmonic lattices respectively, the
integrals in \ref{ari:C.8.6} are
\[
 \int\log|C|\,d\nu=\log2,\qquad
 \int\log|C|\,d\nu=\frac32\log3-2\log2.
\]
They recover $N^2=2$ and $N^4=3$ with the same normalization.
\end{proposition}


\begin{proof}
\begin{enumerate}[label=\arabic*.,leftmargin=*]
\item By homogeneity it suffices to take $C(u)=1-u^m$, with $m=2$
or $m=3$. For $-1<\Re s<1-1/m$ put
\[
 I_m(s)=\int_{\C} |u|^{-1}|1-u^m|^{s-1}dA_u.
\]
The substitution $w=u^m$, with all $m$ sheets counted, gives
\[
 I_m(s)=\frac1m B_{\C}\left(\frac1{2m},\frac{s+1}2\right),
 \quad
 B_{\C}(a,b)=\pi\frac{\Gamma(a)\Gamma(b)\Gamma(1-a-b)}
                   {\Gamma(1-a)\Gamma(1-b)\Gamma(a+b)}.
\]
The complex beta identity \cite[\S\S1.1--1.2, equations~(1.6)--(1.7),
pp.~38--39]{Neretin2024Beta} applies with parameter pairs $a\mid a$ and $b\mid b$.
In that notation $\Gamma^{\mathbb C}(a\mid a)=\Gamma(a)/\Gamma(1-a)$ and
\[
 B_{\C}(a,b)=\int_{\C}|w|^{2a-2}|1-w|^{2b-2}\,dA_w
           =\pi B^{\mathbb C}(a\mid a,b\mid b).
\]
The source's absolute-convergence conditions are
$\Re a>0$, $\Re b>0$ and $\Re(a+b)<1$. Under
$a=1/(2m)$ and $b=(s+1)/2$, these give exactly the stated strip.
The strict inequalities give integrable power bounds at $u=0$, at each
root of $1-u^m$, and at infinity. On a sufficiently small closed
neighbourhood of $s=0$, their positive margins also absorb the
logarithmic factor introduced by differentiation. Thus differentiation
under the integral at $s=0$ is justified.

\item Consequently, with $a=1/(2m)$ and $\psi=\Gamma'/\Gamma$,
\[
 \frac{I_m'(0)}{I_m(0)}
 =\frac12\left(2\psi(1/2)-\psi(1/2-a)-\psi(1/2+a)\right).
\]
The gamma duplication and triplication formulas, differentiated
logarithmically, evaluate this as $\log2$ for $m=2$ and
$\frac32\log3-2\log2$ for $m=3$. These can also be checked
directly from
$\psi(1/2)=-\gamma-2\log2$,
$\psi(1/4)+\psi(3/4)=-2\gamma-6\log2$ and
$\psi(1/3)+\psi(2/3)=-2\gamma-3\log3$.
Equation \eqref{ari:eq:ca-mahler} gives the stated values of $N$.
\end{enumerate}
\end{proof}
